\documentclass[11pt,a4paper]{article}
\usepackage[utf8]{inputenc}
\usepackage[T1]{fontenc}
\usepackage{amsmath,amssymb,amsfonts}
\usepackage{geometry}
\usepackage{booktabs}
\usepackage{enumitem}
\usepackage{listings}
\usepackage{listingsutf8}
\usepackage{xcolor}

\geometry{margin=1in}

\definecolor{primary}{HTML}{1B365D}
\definecolor{secondary}{HTML}{4A777A}
\definecolor{pykeyword}{HTML}{1F4E79}
\definecolor{pystring}{HTML}{A23B00}
\definecolor{pycomment}{HTML}{5C6B73}
\definecolor{pynumber}{HTML}{6C2BD9}

\lstset{
    language=Python,
    basicstyle=\ttfamily\small,
    keywordstyle=\color{pykeyword}\bfseries,
    stringstyle=\color{pystring},
    commentstyle=\color{pycomment}\itshape,
    numberstyle=\tiny\color{pynumber},
    numbers=left,
    numbersep=8pt,
    breaklines=true,
    showstringspaces=false,
    frame=single,
    rulecolor=\color{secondary},
    backgroundcolor=\color{black!2},
    inputencoding=utf8,
    extendedchars=true,
    literate={å}{{\aa}}1 {ä}{{\"a}}1 {ö}{{\"o}}1 {Å}{{\AA}}1 {Ä}{{\"A}}1 {Ö}{{\"O}}1
}

\title{\bfseries\color{primary}Kapitel 1: Kom igång med Python\\ \large Viktiga Begrepp och Studiefrågor (DA2005)}
\author{Studieguide}
\date{\today}

\begin{document}

\maketitle

\section{Kapitel 1: Intro till Python}

\begin{description}[style=nextline, leftmargin=1cm, labelindent=0.5cm]
    \item[\color{primary}Programmering] Att skriva instruktioner (program) som en dator sedan kan köra för att lösa specifika problem.
    \item[\color{primary}Anaconda] En distribution och ett verktygsprogram för Python som används för att hantera paket, bibliotek och virtuella utvecklingen.
    \item[\color{primary}Spyder] En integrerad utvecklingsmiljö (IDE) som följer med Anaconda, utrustad med kodredigerare, syntaxmarkering (syntax highlighting) och en interaktiv konsol.
    \item[\color{primary}Python-tolk (Interpreter)] Det program som läser, analyserar och kör (tolkar) Python-kod rad för rad i realtid.
    \item[\color{primary}Konsol (Console)] Det interaktiva fönstret (t.ex. i Spyder) där man kan skriva enstaka Python-uttryck för att direkt se deras resultat, samt där programmets utskrifter visas.
    \item[\color{primary}Variabel] En namngiven behållare i datorns minne som används för att spara och referera till ett värde.
    \item[\color{primary}Tilldelning (Assignment)] Processen där ett värde lagras i en variabel med hjälp av tilldelningsoperatorn \texttt{=}. Exempel: \texttt{s = 67}.
    \item[\color{primary}Datatyp] En klassificering av data som talar om för tolken hur värdet ska användas. De tre primära typerna i kapitel 1 är:
    \begin{itemize}
        \item \textbf{Heltal (\texttt{int}):} Representerar heltal utan decimaler (t.ex. \texttt{17}, \texttt{-3}).
        \item \textbf{Flyttal (\texttt{float}):} Representerar reella tal med decimalpunkt (t.ex. \texttt{3.14}, \texttt{1.25}).
        \item \textbf{Sträng (\texttt{str}):} En sekvens av tecken som omges av enkelfnuttar (\texttt{'hej'}) eller dubbelfnuttar (\texttt{"hej"}).
    \end{itemize}
    \item[\color{primary}Aritmetiska operationer] Matematiska operationer för beräkningar:
    \begin{itemize}
        \item \textbf{Heltalsdivision (\texttt{//}}): Dividerar två tal och avrundar nedåt till närmaste heltal.
        \item \textbf{Exponentiering (\texttt{**}}): Upphöjt till (t.ex. \texttt{2**3} ger 8).
        \item \textbf{Modulus/Rest (\texttt{\%}}): Returnerar resten vid en heltalsdivision (t.ex. \texttt{5 \% 2} ger 1).
    \end{itemize}
    \item[\color{primary}Reserverade ord (Keywords)] Ord som har en specifik fördefinierad betydelse i Python och som därför inte får användas som variabelnamn (t.ex. \texttt{and}, \texttt{if}, \texttt{for}, \texttt{True}).
    \item[\color{primary}Identifierare (Identifiers)] Namn på resurser i programmet, såsom variabler, funktioner eller klasser. De måste börja med en bokstav eller understreck (\texttt{\_}) och får inte börja med en siffra.
    \item[\color{primary}Litteral] Ett fast, konstant värde som skrivs direkt i programkoden (t.ex. talet \texttt{2.5} eller strängen \texttt{'SU'}). Bör helst inte spridas okontrollerat i koden (undvik s.k. ``hårdkodning'').
\end{description}


\section*{Instuderingsfrågor (Review Questions)}

\begin{enumerate}[label=\bfseries\color{primary}Fråga \arabic*:, leftmargin=1.8cm]
    \item {\bfseries Vad är skillnaden mellan en Python-tolk och en utvecklingsmiljö som Spyder?}
    \\ {\itshape Svar: Tolken kör själva koden, medan Spyder är ett verktyg/redigeringsprogram som gör det lättare för programmeraren att skriva, felsöka och organisera koden.}
    
    \item {\bfseries Varför är ordningsföljden på instruktioner viktig vid tilldelning av variabler? Ge ett exempel på en otillåten ordning.}
    \\ {\itshape Svar: Python körs sekventiellt (rad för rad). Du kan inte använda en variabel i en beräkning innan den har definierats. Till exempel ger följande kod felmeddelandet \texttt{NameError}:}
    \begin{verbatim}
    v = s / t  # Fel! s och t är inte definierade än.
    s = 67
    t = 1.25
    \end{verbatim}

    \item {\bfseries Förklara skillnaden i resultat mellan operatorerna \texttt{/}, \texttt{//}, och \texttt{\%} med hjälp av talen 7 och 3.}
    \\ {\itshape Svar:}
    \begin{itemize}
        \item \texttt{7 / 3} ger flyttalet \texttt{2.3333333333333335} (vanlig division).
        \item \texttt{7 // 3} ger heltalet \texttt{2} (heltalsdivision, kapar decimalerna).
        \item \texttt{7 \% 3} ger resten \texttt{1} (modulus, eftersom $7 = 2 \times 3 + 1$).
    \end{itemize}

    \item {\bfseries Vad händer om du kör koden \texttt{'hej' * 3} respektive \texttt{'hej' + 3}? Förklara varför.}
    \\ {\itshape Svar:} 
    \begin{itemize}
        \item \texttt{'hej' * 3} duplicerar strängen och ger \texttt{'hejhejhej'}. Detta är en tillåten operation (sträng multiplicerat med heltal).
        \item \texttt{'hej' + 3} ger ett \texttt{TypeError} eftersom man inte kan addera (konkatenera) en sträng med ett heltal (\texttt{int}) direkt utan att först konvertera talet till en sträng.
    \end{itemize}

    \item {\bfseries Vilka av följande är giltiga identifierare (variabelnamn) i Python? Om de är ogiltiga, förklara varför.}
    \begin{enumerate}[label=\alph*)]
        \item \texttt{my\_variable}
        \item \texttt{2nd\_value}
        \item \texttt{class}
        \item \texttt{Value\_17}
    \end{enumerate}
    {\itshape Svar:}
    \begin{itemize}
        \item a) och d) är {\bfseries giltiga}.
        \item b) är {\bfseries ogiltig} eftersom en identifierare inte får börja med en siffra.
        \item c) är {\bfseries ogiltig} eftersom \texttt{class} är ett reserverat ord (keyword) i Python.
    \end{itemize}

    \item {\bfseries Vad menas med att ``hårdkoda'' ett värde, och varför bör man undvika litteraler spridda i koden enligt kompendiet?}
    \\ {\itshape Svar: Att hårdkoda innebär att skriva in konstanta värden (litteraler) direkt i beräkningar spridda över programmet. Det bör undvikas eftersom det försvårar uppdateringar (om värdet ändras måste man leta upp alla ställen manuellt) och ökar risken för programmeringsfel. Det är bättre att definiera värdet en gång i en tydligt namngiven variabel längst upp i koden.}
\end{enumerate}

\section{Kapitel 2: intro till python}

\begin{description}[style=nextline, leftmargin=1cm, labelindent=0.5cm]
    \item[\color{primary}Syntax] Regler för hur kod skrivs för att Python ska förstå den.
    \item[\color{primary}Semantik] Innebörden eller betydelsen av vad koden gör.
    \item[\color{primary}Datatyper] Python har inbyggda typer som \texttt{str}, \texttt{int}, \texttt{float}, \texttt{complex}, \texttt{bool} och \texttt{NoneType}.
    \item[\color{primary}Typomvandling] Kan vara implicit (automatisk) eller explicit (t.ex. \texttt{int(x)}).
    \item[\color{primary}Kortslutning (Short circuiting)] När logiska operatorer (\texttt{and}/\texttt{or}) avbryter utvärderingen så fort resultatet är garanterat.
    \item[\color{primary}Modulär kod] Program uppdelade i mindre, fristående funktioner.
    \item[\color{primary}Docstrings] Dokumentationssträngar knutna till funktioner, läsbara via \texttt{help()}.
    \item[\color{primary}Scope] Räckvidden för en variabel (global eller lokal).
\end{description}

\begin{enumerate}[label=\bfseries\color{primary}Fråga \arabic*:, leftmargin=1.8cm]
    \item \textbf{Vad är skillnaden mellan SyntaxError och semantiska fel?}
    \\ \textit{Svar:} Syntaxfel innebär att koden inte följer språkets grammatik, vilket gör att den inte kan köras. Semantiska fel innebär att koden är korrekt skriven men utför fel logik, vilket resulterar i oväntade utdata.
    
    \item \textbf{Varför representeras flyttal som approximationer?}
    \\ \textit{Svar:} Python representerar flyttal genom bråk i bas 2. Det är matematiskt omöjligt att representera alla reella tal exakt med ett begränsat antal bitar på en dator.
    
    \item \textbf{Hur ändrar man en global variabel inuti en funktion?}
    \\ \textit{Svar:} Genom att använda det reserverade ordet \texttt{global} inuti funktionen deklarerar man att man avser att ändra den globala variabeln snarare än att skapa en lokal variabel med samma namn.
    
    \item \textbf{Vad händer om en funktion saknar \texttt{return}?}
    \\ \textit{Svar:} Om \texttt{return} saknas returnerar funktionen värdet \texttt{None} som standard.
    
    \item \textbf{Varför är \texttt{elif} användbart?}
    \\ \textit{Svar:} Det förbättrar kodens läsbarhet genom att ersätta krångliga, nästlade \texttt{if-else}-strukturer, vilket gör logiken tydligare att följa.
\end{enumerate}

\section{Kapitel 3: Listor och iteration}

\begin{description}[style=nextline, leftmargin=1cm, labelindent=0.5cm]
    \item[\color{primary}Lista] En sekvens av element (av valfri typ) inneslutna i hakparenteser \texttt{[]}.
    \item[\color{primary}Iteration] Processen att automatisera repetitiva uppgifter.
    \item[\color{primary}for-loop] Används för att iterera över element i en sekvens.
    \item[\color{primary}while-loop] Används för att repetera kod så länge ett villkor är sant.
\end{description}

\begin{enumerate}[label=\bfseries\color{primary}Fråga \arabic*:, leftmargin=1.8cm]
    \item \textbf{Varför är listor viktiga?}
    \\ \textit{Svar:} De gör det möjligt att samla och hantera flera element effektivt, vilket underlättar beräkningar och databehandling.
    
    \item \textbf{Vad är skillnaden mellan for och while?}
    \\ \textit{Svar:} \texttt{for}-loopar används för att iterera över samlingar, medan \texttt{while}-loopar repeterar kod baserat på ett sant villkor.
\end{enumerate}

\section{Kapitel 4: Filer och dictionaries}

\begin{description}[style=nextline, leftmargin=1cm, labelindent=0.5cm]
    \item[\color{primary}Datastruktur] En struktur för att organisera och samla data logiskt.
    \item[\color{primary}Uppslagstabell] En datastruktur som kan indexeras med andra typer än bara heltal.
    \item[\color{primary}Filhantering] Tekniker för att läsa från och skriva till filer på disk.
\end{description}

\section*{Besvarade Instuderingsfrågor}
\begin{enumerate}[label=\bfseries\color{primary}Fråga \arabic*:, leftmargin=1.8cm]
    \item \textbf{Vad skiljer en uppslagstabell från en lista?}
    \\ \textit{Svar:} Listor indexeras via heltal, medan uppslagstabeller tillåter indexering via andra datatyper.
    
    \item \textbf{Vad är syftet med filhantering?}
    \\ \textit{Svar:} Det möjliggör permanent lagring av data samt inläsning av extern data som inte finns direkt i programkoden.
\end{enumerate}

\section{Kapitel 5: Felhantering och undantag}

\begin{description}[style=nextline, leftmargin=1cm, labelindent=0.5cm]
    \item[\color{primary}Felhantering] Metoder för att hantera programfel (t.ex. krascher).
    \item[\color{primary}Särfall (Exception)] Signalering av fel som uppstått under körning.
    \item[\color{primary}try/except] Pythons verktyg för att fånga och hantera uppkomna fel.
    \item[\color{primary}Listomfattning] Kompakt syntax för att skapa och bearbeta listor.
\end{description}

\begin{enumerate}[label=\bfseries\color{primary}Fråga \arabic*:, leftmargin=1.8cm]
    \item \textbf{Varför är särfall bättre än felkoder?}
    \\ \textit{Svar:} Felkoder gör koden svårläst och komplex p.g.a. ständiga kontroller. Särfall separerar felhantering från den ordinarie logiken.
    
    \item \textbf{Vad är fördelen med listomfattning?}
    \\ \textit{Svar:} Det möjliggör mer kompakt och effektiv kod vid skapande av listor jämfört med traditionella loopar.
\end{enumerate}

\section{Kapitel 6: Sekvenser och generatorer}

\begin{description}[style=nextline, leftmargin=1cm, labelindent=0.5cm]
    \item[\color{primary}Sekvenstyp] Datatyp med gemensamma operationer (t.ex. indexering, slice).
    \item[\color{primary}Muterbar] Objekt som kan ändras efter skapande (t.ex. listor).
    \item[\color{primary}Icke-muterbar] Objekt som är konstanta (t.ex. tupler, strängar).
\end{description}

\begin{enumerate}[label=\bfseries\color{primary}Fråga \arabic*:, leftmargin=1.8cm]
    \item \textbf{Vad är skillnaden mellan muterbara och icke-muterbara objekt?}
    \\ \textit{Svar:} Muterbara objekt kan ändras efter skapande, medan icke-muterbara objekt är konstanta.
    
    \item \textbf{Varför är muterbarhet viktigt?}
    \\ \textit{Svar:} Det påverkar hur objekt hanteras i funktioner och kan leda till oväntade sidoeffekter om man inte är medveten om att originalobjektet ändras.
\end{enumerate}

\section{Kapitel 7: Moduler, biblotek och programstruktur}

\begin{description}[style=nextline, leftmargin=1cm, labelindent=0.5cm]
    \item[\color{primary}Sekvenstyp] Datatyp med gemensamma operationer (t.ex. indexering, slice).
    \item[\color{primary}Muterbar] Objekt som kan ändras efter skapande (t.ex. listor).
    \item[\color{primary}Icke-muterbar] Objekt som är konstanta (t.ex. tupler, strängar).
\end{description}

\begin{enumerate}[label=\bfseries\color{primary}Fråga \arabic*:, leftmargin=1.8cm]
    \item \textbf{Vad är skillnaden mellan muterbara och icke-muterbara objekt?}
    \\ \textit{Svar:} Muterbara objekt kan ändras efter skapande, medan icke-muterbara objekt är konstanta.
    
    \item \textbf{Varför är muterbarhet viktigt?}
    \\ \textit{Svar:} Det påverkar hur objekt hanteras i funktioner och kan leda till oväntade sidoeffekter om man inte är medveten om att originalobjektet ändras.
\end{enumerate}

\begin{lstlisting}
    '''
    short/modular (variables, functions, files) 
    variables (nouns,repeating stuff should be global constants)
    functions (action on what, return for local > global=constant, keep be math alone, don't return int/bool/fixed stuff/unspecified numbers, several returns > one)
    documentation (what, arg & returns - name, type, description)
    OOP (code skeleton)
    Tests
    ! main
    commenting (why you chose a method, use this when) - no syntax explinations, no examples, no long sentances

    Type hinting (def(arg: type) -> return_type)
    '''
\end{lstlisting}

\section{Kapitel 8: Funktionell programmering}

\begin{description}[style=nextline, leftmargin=1cm, labelindent=0.5cm]
    \item[\color{primary}Funktionell programmering] bygger på att man anropar, sätter ihop och modifierar funktioner.
    \item[\color{primary}Rekursion] En teknik där en funktion anropar sig själv för att lösa ett problem.
    \item[\color{primary}Högre ordningens funktioner] Funktioner som tar andra funktioner som argument eller returnerar funktioner.
    \item[\color{primary}Anonyma funktioner (lambda)] Funktioner utan namn som definieras med \texttt{lambda}-syntaxen. Används ofta som argument till högre ordningens funktioner då de inte behöver återanvändas.
    \item[\color{primary}Svansrekursion] En form av rekursion där det rekursiva anropet är det sista som sker i funktionen. Vissa språk kan optimera detta till en loop för att spara minne.
\end{description}

\begin{enumerate}[label=\bfseries\color{primary}Fråga \arabic*:, leftmargin=1.8cm]
    \item \textbf{Vad innebär rekursion?}
    \\ \textit{Svar:} Rekursion är när en funktion anropar sig själv. Det är ett alternativ till loopar och kan göra koden mer elegant.
    
    \item \textbf{Vad kännetecknar funktionell programmering?}
    \\ \textit{Svar:} Det bygger på att använda, kombinera och modifiera funktioner snarare än att styra programflödet med loopar och variabler som ändras.

    \item \textbf{Varför behöver rekursion ett basfall?}
    \\ \textit{Svar:} Ett basfall behövs för att stoppa rekursionen och förhindra oändliga anrop.

\end{enumerate}

\begin{lstlisting}[language=Python]
    def pascal(n):
    '''
    Ex:
    1. pascal(3) exekveras:
       n = 3 (inte 1, gar vidare)
       Anropar pascal(2) for att fa 'prev_row'
    2. pascal(2) exekveras:
       n = 2 (inte 1, gar vidare)
       Anropar pascal(1) for att fa 'prev_row'
    3. pascal(1) exekveras:
       n = 1 (basfallet uppnatt!)
       RETURVARDE FRAN pascal(1): [1]

    [RETURFASEN - Vagen upp och berakning av varden]
    4. Tillbaka i pascal(2):
        Mottagit 'prev_row' = [1]
        Initierar 'new_elements' = []
        Beraknar loopen: range(len([1]) - 1) -> range(0) -> Loopen kors 0 ganger
        Satter ihop raden: [1] + [] + [1]
        RETURVARDE FRAN pascal(2): [1, 1]
    5. Tillbaka i pascal(3):
        Mottagit 'prev_row' = [1, 1]
        Initierar 'new_elements' = []
        Beraknar loopen: range(len([1, 1]) - 1) -> range(1) -> Loopen kors for i = 0
            * i = 0: prev_row[0] + prev_row[1] -> 1 + 1 = 2
            * new_elements.append(2) -> 'new_elements' ar nu [2]
        Satter ihop raden: [1] + [2] + [1]
        RETURVARDE FRAN pascal(3): [1, 2, 1]
    '''
        if n == 1:
            return [1]
        
        prev_row = pascal(n - 1)
        new_elements = []
        
        for i in range(len(prev_row) - 1):
            new_elements.append(prev_row[i] + prev_row[i+1])
            
        return [1] + new_elements + [1]
\end{lstlisting}

\section{Kapitel 9: OOP och klasser}

\begin{description}[style=nextline, leftmargin=1cm, labelindent=0.5cm]
    \item[\color{primary}Objektorientering] är ett programmeringsparadigm eller programmeringsstil där man samlar data och datastrukturer med de funktioner som kan tillämpas på dem.
    \item[\color{primary}Klass, objekt, instans] En klass är en mall för att skapa objekt. Ett objekt är en instans av en klass (dvs. ett konkret exempel på klassen). Skillnaden mellan objekt och instans är att objekt är en mer generell term, medan instans syftar på ett specifikt objekt som skapats från en klass. 
    \item[\color{primary}Attribut] En variabel som är knuten till ett specifikt objekt och beskriver dess nuvarande tillstånd. De är som objektets variabler. Det finns instansattribut (knutna till ett specifikt objekt) och klassattribut (delade av alla objekt i klassen).
    \item[\color{primary}Metod] En funktion som är knuten till ett specifikt objekt och kan manipulera dess attribut eller utföra operationer relaterade till objektet. De är som funktioner till objektet.
    \item[\color{primary}Konstruktor] En speciell metod som används för att skapa och initiera objekt av en klass. I Python är konstruktorn metoden \texttt{\_\_init\_\_()}. Så när man anropar klassen igen för att skapa ett nytt objekt, körs konstruktorn automatiskt för att sätta upp objektets initiala tillstånd.
    \item[\color{primary}Metod] En funktion som är knuten till ett specifikt objekt.
    \item[\color{primary}Punktnotation] Syntaxen \texttt{objekt.metod()} för att anropa metoder.
    \item[\color{primary}Self] En referens till det aktuella klassensobjektet, används för att komma åt objektets attribut och metoder. Kan kallas annat men self är konvention.
    \item[\color{primary}Privata och publika attribut/metoder] Publika attribut/metoder kan nås utanför klassen, medan privata är avsedda att endast användas inom klassen. I Python markeras privata med ett understreck (\_) i början av namnet.
    \item[\color{primary}Modulär programmering] Att dela upp program i mindre, fristående enheter (moduler) som kan återanvändas och underhållas separat vilket leder till mindre kod att ändra vid behov.
\end{description}

\begin{lstlisting}[language=Python]
    # Syntax for classes
    class KlassNamn: # Bör namnges med stor bokstav i början, va ett substantiv tex Bilar

        competition = "Formula 1" # Klassattribut, delad av alla objekt i klassen

        def __init__(self, arg1, arg2): # Instansatribut, unik för varje objekt genom att data läggs till
            self.attribut1 = arg1 # Här skapas attribut/variabler tex ålder, vikt
            self.attribut2 = arg2
        
        def metod(self, attribut1, attribut2): # funktion/metod som beräknar bilens värde
            värde = self.attribut1 * 1000 + self.attribut2 * 500 # Obs även enkla funktioner som len() behövs läggas till så de funkar på din klass
            return värde

        # Lägg till enkla funktioner vi vill ha
        def __str__(self): # __str__() är en speciell metod som definierar hur objektet ska representeras som en sträng
            return f"Objekt med attribut1: {self.attribut1}, attribut2: {self.attribut2}, competition: {KlassNamn.competition}"
    
    # Syntax för att skapa ett objekt (instans) av klassen
    objekt = KlassNamn(ålder1, vikt2) # objekt kan nu manipulera hela datan från Bilar med hjälp av enkel punktnotation
    objekt.attribut3 = värde3 # Här kan man lägga till fler attribut, obs kan skrivas över
    objekt.metod() # Här kan man anropa metoder som redan finns i klassen
    objekt2 = KlassNamn(ålder2, vikt3) # Här skapas ett nytt objekt med annan data som en kopia
    print(objekt2.attribut2) # För att använda funktioner som redan finns på objektet, kalla på dess info/attribut
\end{lstlisting}

\begin{lstlisting}[language=Python]
    # Privata och publika attribut/metoder
    class Bil:
        def __init__(self, ålder, vikt):
            self._ålder = ålder  # Privata attribut (bör inte ändras utanför klassen)
            self._vikt = vikt    # Privata attribut (bör inte ändras utanför klassen)

        def beräkna_värde(self):
            return self._ålder * 1000 + self._vikt * 500  # Publik metod som använder privata attribut

    min_bil = Bil(5, 1500) # man kan manipulera med beräkna_värde() men inte ändra _ålder eller _vikt direkt om man inte skriver...
    min_bil._ålder = 10  # ...så här, men det är inte rekommenderat eftersom det bryter mot principen om inkapsling.
\end{lstlisting}

\begin{lstlisting}[language=Python]
    # Sets
    a = {1, 2, 3, 4, 5} # Vilket är samma som {1, 1, 2, 3, 4, 5, 5}
    b = {4, 5, 6, 7, 8}
    print(a | b)  # Union
    print(a & b)  # Intersection
    print(a - b)  # Difference
    print(a ^ b)  # Symmetric difference = (a - b) | (b - a)
\end{lstlisting}

\begin{enumerate}[label=\bfseries\color{primary}Fråga \arabic*:, leftmargin=1.8cm]
    \item \textbf{Varför använda objektorientering?}
    \\ \textit{Svar:} Knyta samman mycket data som relaterar till varandra (hör till samma objekt) vilket ger bra struktur med relaterad funktionalitet på samma ställe. Också att kunna definiera egna datatyper som hjälper abstraktion genom att kunna gömma detaljer.
\end{enumerate}

\section{Kapitel 10: OOP arv}
\begin{description}[style=nextline, leftmargin=1cm, labelindent=0.5cm]
    \item[\color{primary}Arv] Mekanism för att återanvända och utöka kod genom subklasser.
    \item[\color{primary}Basklass] Den ursprungliga klassen som funktionalitet ärvs ifrån.
    \item[\color{primary}Subklass] Den klass som ärver från en basklass.
\end{description}

\begin{enumerate}[label=\bfseries\color{primary}Fråga \arabic*:, leftmargin=1.8cm]
    \item \textbf{Varför använder vi arv?}
    \\ \textit{Svar:} För att återanvända kod på ett strukturerat sätt och för att kunna bygga vidare på existerande klasser.
    
    \item \textbf{Hur definieras en subklass i Python?}
    \\ \textit{Svar:} Man anger basklassen inom parentes i subklassens definition, t.ex. \texttt{class Sub(Bas):}.
\end{enumerate}

\section{Kapitel 11: OOP arv mer}

\begin{description}[style=nextline, leftmargin=1cm, labelindent=0.5cm]
    \item[\color{primary}Multipelt arv] När en klass ärver från två eller fler superklasser.
    \item[\color{primary}Diamantproblemet] En tvetydighet i arvshierarkin som kan uppstå vid multipelt arv.
\end{description}

\begin{enumerate}[label=\bfseries\color{primary}Fråga \arabic*:, leftmargin=1.8cm]
    \item \textbf{Vad är diamantproblemet?}
    \\ \textit{Svar:} Det är en konflikt i arvshierarkin där en subklass ärver från två klasser som delar samma basklass, vilket skapar osäkerhet om vilken metod som ska anropas.
    
    \item \textbf{Varför är multipelt arv kontroversiellt?}
    \\ \textit{Svar:} Eftersom det kan göra arvshierarkin oöverskådlig och skapa konflikter som inte finns i enkla arvsstrukturer.
\end{enumerate}

\section{Kapitel 12: Defensiv programmering}

\begin{description}[style=nextline, leftmargin=1cm, labelindent=0.5cm]
    \item[\color{primary}Defensiv programmering] Ansats för att identifiera och hantera fel tidigt.
    \item[\color{primary}Assert] Kontroll av logiska antaganden i koden.
    \item[\color{primary}unittest] Modul för att automatisera och systematisera kodtestning.
\end{description}

\begin{enumerate}[label=\bfseries\color{primary}Fråga \arabic*:, leftmargin=1.8cm]
    \item \textbf{Varför använda assert?}
    \\ \textit{Svar:} För att verifiera att programmets tillstånd stämmer och fånga logiska fel tidigt.
    
    \item \textbf{Vad gör unittest-modulen?}
    \\ \textit{Svar:} Den tillhandahåller ett ramverk för att organisera och köra automatiserade tester systematiskt.
\end{enumerate}

\subsection{Test kap. 1-4}

\textbf{Kom ihåg!!!}
\begin{itemize}
    \item Stil: En förstående stil. Engelska, förklarande namn ($mean_vel$ vs v vs $mean_velocity$), syntax/semantik
    \item Reserverade ord: and, if, for, True, False, None
    \item Identifierare: Börja med bokstav eller understreck, inte siffra
    \item Hårdkodning: Undvik litteraler spridda i koden, definiera en gång i en tydligt namngiven variabel
\end{itemize}


\textbf{Teorifrågor}:
\begin{enumerate}
    \item Vad är en funktions kropp, huvud, parametrar, argument, dokumentsträng
    \item Vad är skillnaden på imperativ, funktionell, OOP programmering
    \item När används pass vs continue: pass används som en platshållare för framtida kod, medan continue används för att hoppa över resten av koden i en loop och fortsätta med nästa iteration.
    \item När används for vs while loopar: for används när man vet hur många iterationer
    \item När används .copy() tex för variabler och listor: .copy() används för att skapa en ny kopia av en lista eller ett objekt, vilket gör att ändringar i den nya kopian inte påverkar originalet.
    \item När ska man använda en lista vs en dict: array vs hash-tabell om man vill nå top och botten eller allt
    \item Vad är hårdisk och operativsystem: hårddisken är en lagringsenhet som används för att permanent lagra data, medan operativsystemet är den programvara tex Windows, Linux, macOS
    \item Varför ska man f.close() filer: För att frigöra resurser och säkerställa att alla data skrivs till filen korrekt.
    \item När används filerna sys.stdin out derr: tangetnbord, konsol, felmeddelanden
    \item Saker jag hoppat över: bra praxis för att programera tillsammans 4.4.3 samt representation av data 4.5
\end{enumerate}


\textbf{Live Coding}:
\begin{enumerate}
    \item Skapa en python fil. Kör filen samt förbättra denna kod:
    \begin{lstlisting}
    # Denna kod är till för att söka igenom en lista.
    # Den kollar varje element i sekvensen ett i taget för att hitta ett specifikt värde.
    # Om elementet finns i listan returneras positionen där det hittades.

    target = 10

    def search_element(items, target):
        global target
        """
        Koden går igenom hela listan och jämför elementen.
        Om rätt värde hittas sparas indexet.
        """
        found_index = -1
        for i in range(len(items)):
            if items[i] == target:
                found_index = i
                pass
                break
            else:
                pass
            pass
        return found_index

    # Anrop
    resultat = search_element([3, 8, 10, 15])
    \end{lstlisting}
    \item Hitta alla fel i denna kod:
    \begin{lstlisting}
        räknare = 0

        def skapa_skylt_rapport():
            ägare = "Karolina'
            hälsning = "Skylt skapad av" + "Karolina"
            dekorationskant = 2 * "="
            area = 0^2 * 3.14
            
            if 0.1 + 0.2 == 0.3:
                print("Precisionen är exakt!")
                
            if bool(0) = True: 
                print("Status: Giltig radie")
                
            if area => 50:
                status_kod = 10
                if status_kod = 10:
                    rapport_text = hälsning + " - Stor skylt"
                    
            print(rapport_text)
            räknare = räknare + 1

        skapa_skylt_rapport()
    \end{lstlisting}
    \item Förklara denna kod:
    \begin{lstlisting}
        x, y = 2, 3

        k = x // y
        r = x % y

        text = f"Resultat:\nKvot: {k}\nRest: {r}"
        f = float(k)
        c = 2 + 3j
        l = x < y
        t = None

        print(text)
        print(type(text), type(x), type(f), type(c), type(l), type(t))

        help(int)
        help(None)
    \end{lstlisting}
    \item Brev till framtida jag
        Skapa ett Python-program som i en loop samlar in användardata (namn, ålder, intressen) och lagrar den i en dictionary för att generera ett "framtidsbrev". 
        

        Programmet beräknar åldern om 5 år, avgör om nuvarande ålder är jämn eller udda, samt hanterar negativa åldrar med felmeddelanden. Om åldern överstiger 100 år avbryts loopen (break) och ett meddelande om att användaren är en "Du är en legendarisk drake och vet redan allt" skrivs ut.


        Slutligen skrivs en sammanfattning ("framtidsbrevet") ut i konsolen baserat på dictionary-datan. Loopen körs tills användaren anger "exit" eller avbryter med Ctrl + C (KeyboardInterrupt).
    \item I denna uppgift ska du arbeta med listan 'data' som skapats i förberedelsekoden. 
    Utför följande operationer i ordning:


    1. Ta fram och skriv ut det näst sista värdet.

    2. Lägg till ett nytt värde i slutet av listan.

    3. Ta bort det 3:e elementet (index 2) med 'del'.
    
    4. Ta bort det 3:e elementet (index 2) med 'pop()' och skriv ut det returnerade värdet.
    
    5. Ta bort alla förekomster av "joker" från listan med 'remove()'.
    
    6. Byt ut det 3:e elementet (index 2) mot ett annat valfritt värde.
    
    7. Skriv ut elementen från index 3 till 9, men bara varannat element.
    
    8. Skriv ut listan baklänges med slicing (steglängd -1).
    
    9. Skriv ut alla element i listan med en 'for'-loop.
    
    10. Skriv ut alla element i listan med en 'while'-loop.
    
    11. Skriv ut hur lång listan är (antal element).
    
    12. Kontrollera om talet 7 finns i listan och skriv ut ett booleskt resultat.
    
    13. Ta fram och skriv ut vilket index "katt" har.
    
    14. Räkna och skriv ut hur många förekomster av "katt" som finns i listan.
    
    15. Vänd på listan baklänges direkt på plats med metoden '.reverse()' och skriv ut den.
    
    16. Lägg till strängen "hund" längst bak i listan med metoden '.insert()'.


    Här är koden du behöver:
    \begin{lstlisting}
        # Förberedelsekod: Skapar startlistan med tillräckligt med element
    data = [
        "apa",
        "björn",
        "joker",
        "katt",
        7,
        "joker",
        "elefant",
        "katt",
        "fisk",
        "gepardt",
        "häst",
        "joker",
        10,
    ]

    print("Ursprunglig lista:", data)
    print("-" * 50)
    \end{lstlisting}

    \item Skapa en funktion \texttt{generera\_x(start, stop, step=1)} som returnerar en lista med x-värden med \texttt{list(range(start, stop, step))}. Paramentern 'step' ska ha standardvärdet 1. Beräkna tillhörande y-värden (t.ex. y = $x^2$) och rita grafen med matplotlib.pyplot.
    \item Givet meningen: "jag grät när jag inte fick mina björnbär" och de förbjudna bokstäverna: 'åäö'. Skriv ett Python-program som:
            
        1. Kontrollerar om meningen innehåller någon av de förbjudna bokstäverna.
        
        2. Identifierar och skriver ut alla enskilda ord i meningen som innehåller någon av dessa bokstäver.
    \item Skapa och manipulera en dictionary (uppslagstabell) i Python genom att utfora foljande steg:
        1. Skapa en dictionary som representerar en frugtlager-status med nycklar (fruktnamn) och varden (antal i lager).
        2. Hamta och skriv ut antalet "applen" med hjalp av metoden .get().
        3. Uppdatera lagersaldot for "banan" till ett nytt varde.
        4. Iterera igenom hela din dictionary med metoden .items() och skriv ut varje nyckel och varde pa ett snyggt satt.
    \item I denna uppgift ska du använda 'with open(...) as fil' för att hantera filer utan att behöva stänga dem manuellt. Börja med att läsa in en textfil med ett kärleksbrev till katter genom att använda metoden file.readlines(). Därefter ska du skriva om koden så att alla förekomster av ordet katt ersätts med hund, och spara det nya brevet i en helt ny fil. Avslutningsvis ska du läsa in den nya filen och skriva ut hela dess innehåll i konsolen.
    
    Här är förberdelsekod du behöver:
    \begin{lstlisting}
        # Förberedelsekod: Skapar ursprungsposten
        meddelande = """Kära katt, du är den underbaraste katt som finns.
        När en katt spinner blir man alltid glad.
        Tack för att du är min katt!"""

        with open("karleksbrev_katt.txt", "w") as fil:
            fil.write(meddelande)
    \end{lstlisting}
    \item Skapa en 3x3-matris [[1, 2, 3], [4, 5, 6], [7, 8, 9]] med nästlade for-loopar. Skriv ut matrisen rad för rad med en nästlad loop, samt beräkna och skriv ut summan av alla element.
    \item Gör en while loop som printar nummer och räknar ner från ner från 5 till 0
    \item Skriv en funktion som printar odd eller even vad en integer är som tas som input
    \item Skriv en funktion som ger boolska värden om ett ord är längre än 8 tecken eller inte
\end{enumerate}


Facit för live coding:
\begin{enumerate}
    \item Problem 1
    \begin{lstlisting}
        """Algoritm: Lineär sökning"""  # FIX: Bytte från # till docstring, lade till algoritmnamn och tog bort förklaring av koden

        # FIX: Tagit bort den globala variabeln target = 10 (gjorts om till lokal/parameter)

        def search_element(items, target=10):  # FIX: Använder standardvärde target=10 i parametern
            # FIX: Tagit bort 'global target' och hanterar target lokalt via parametern
            found_index = -1
            for i in range(len(items)):
                if items[i] == target:
                    found_index = i
                    # FIX: Tog bort onödig 'pass'
                    break
                # FIX: Tog bort onödig 'else:' och 'pass'
            # FIX: Tog bort onödig 'pass'
            return found_index

        # Anrop
        resultat = search_element([3, 8, 10, 15])  # Använder standardvärdet 10 automatiskt
    \end{lstlisting}
    \item Problem 2
    \begin{lstlisting}
        räknare = 0  # Global variabel

        def skapa_skylt_rapport():
            global räknare  # FIX: Lade till 'global' för att tillåta ändring av den globala variabeln inuti funktionen
            
            ägare = "Karolina"  # FIX: Ändrade enstaka fnutt (') till dubbelfnutt (") för att matcha strängens start
            hälsning = "Skylt skapad av " + ägare  # FIX: Lade till mellanrum i strängen samt använde variabeln 'ägare'
            dekorationskant = "=" * 2  # FIX: Ändrade till att multiplicera tecknet med antalet gånger (renare stil/semantik)
            
            radie = 3  # FIX: Bytte ut XOR-operatorn (^) mot potensoratorn (**) nedan 
            area = radie**2 * 3.14  # FIX: Använder variabeln radie istället för hårdkodad 0
            
            if abs((0.1 + 0.2) - 0.3) < 1e-9:  # FIX: Bytte exakt likhet (==) mot en toleransjämförelse p.g.a. flyttals- / maskinprecision
                print("Precisionen är exakt!")
                
            if not bool(0):  # FIX: Ändrade = till == (eller 'not bool(0)' eller != True) då bool(0) utvärderas till False
                print("Status: Giltig radie")
                
            rapport_text = hälsning + " - Standard skylt"  # FIX: Deklarerar variabeln utanför if-blocket så den inte får för kort räckvidd (ScopeError)
            
            if area >= 50:  # FIX: Ändrade den felaktiga operatorn => till den korrekta >= (större än eller lika med)
                status_kod = 10
                if status_kod == 10:  # FIX: Bytte tilldelningsoperatorn (=) mot jämförelseoperatorn (==) inuti villkoret
                    rapport_text = hälsning + " - Stor skylt"
                    
            print(rapport_text)
            räknare = räknare + 1 

            '''
            I detta fall kommer output bli:
            Skylt skapad av Karolina - Standard skylt

            Om area >= 50, output blir:
            Skylt skapad av Karolina - Stor skylt
            '''

        skapa_skylt_rapport()
    \end{lstlisting}
    \item Problem 3
    \begin{lstlisting}
        # 1. Parallell tilldelning av int-värden
        x, y = 2, 3

        # 2. Heltalsdivision (//) och Modulo/Rest (%)
        kvot = x // y  # Resultat: 0
        rest = x % y   # Resultat: 2 ty 2 = 0 * 3 + 2

        # 3. Olika datatyper i Python
        text = f"Resultat:\nKvot: {kvot}\nRest: {rest}"  # str med radbrytning (\n)
        flyttal = float(kvot)                           # float (0.0)
        komplext_tal = 2 + 3j                           # complex
        är_mindre = x < y                                # bool (True)
        inget_värde = None                              # NoneType (None)

        # Skriv ut strängen
        print(text)

        # 4. Undersök typer och dokumentation med help() i terminalen
        print("\n--- Datatyper ---")
        print(type(text), type(x), type(flyttal), type(komplext_tal), type(är_mindre), type(inget_värde))
        # Output: <class 'str'> <class 'int'> <class 'float'> <class 'complex'> <class 'bool'> <class 'NoneType'>

        # Använd help() för att slå upp en typ eller variabel
        help(int)
        help(None)
    \end{lstlisting}
    \item Problem 4
    \begin{lstlisting}
    while True:
        name = input("Namn (skriv 'exit' för att avsluta): ")
        if name.lower() == "exit":
            break

        age_input = input("Ålder: ")
        if age_input.lower() == "exit":
            break

        # Kontrollera om inmatningen är siffror (eller ett negativt tal)
        if age_input.startswith("-") and age_input[1:].isdigit():
            age = int(age_input)
        elif age_input.isdigit():
            age = int(age_input)
        else:
            print("Ogiltig inmatning! Vänligen ange ålder med siffror.")
            continue

        # Logiska villkor för ålder
        if age < 0:
            print(
                """
    --------------------------------------------------
    FELMEDDELANDE:
    Åldern kan inte vara negativ! Vänligen försök igen.
    --------------------------------------------------
    """
            )
            continue
        elif age > 100:
            print(
                """
    ==================================================
    Du är en legendarisk drake och vet redan allt!
    Programmet avslutas nu.
    ==================================================
    """
            )
            break

        interests = input("Intressen: ")
        if interests.lower() == "exit":
            break

        # Lagra informationen i en dictionary
        user_data = {
            "namn": name,
            "alder": age,
            "framtida_alder": age + 5,
            "jamn_eller_udda": "jämnt" if age % 2 == 0 else "udda",
            "intressen": interests,
        }

        # Utskrift av framtidsbrevet med en flerradig sträng
        framtidsbrev = f"""
    ==================================================
                    FRAMTIDSBREV
    ==================================================
    Hej {user_data['namn']}!

    Du är idag {user_data['alder']} år gammal, vilket är ett {user_data['jamn_eller_udda']} tal.
    Om 5 år kommer du att vara {user_data['framtida_alder']} år gammal.

    Dina registrerade intressen:
    - {user_data['intressen']}
    ==================================================
    """
        print(framtidsbrev)
    \end{lstlisting}
    \item Problem 5
    \begin{lstlisting}
    # Förberedelsekod
    data = [
        "apa",
        "björn",
        "joker",
        "katt",
        7,
        "joker",
        "elefant",
        "katt",
        "fisk",
        "gepardt",
        "häst",
        "joker",
        10,
    ]

    # 1. Näst sista värdet
    print("1. Näst sista värdet:", data[-2])

    # 2. Lägg till ett värde i slutet (append)
    data.append("orm")
    print("2. Efter tillägg av 'orm':", data)

    # 3. Ta bort det 3:e värdet (index 2) med del
    del data[2]
    print("3. Efter del data[2]:", data)

    # 4. Ta bort det 3:e värdet (index 2) med pop och skriv ut returvärdet
    borttaget_varde = data.pop(2)
    print(f"4. Borttaget värde med pop(2): {borttaget_varde}")
    print("   Listan efter pop:", data)

    # 5. Ta bort alla 'joker' med remove
    while "joker" in data:
        data.remove("joker")
    print("5. Efter att alla 'joker' tagits bort:", data)

    # 6. Byt ut det 3:e värdet (index 2) mot ett annat värde
    data[2] = "TIGER"
    print("6. Efter byte av 3:e värdet till 'TIGER':", data)

    # 7. Print index 3 till 9 men bara varannan (slicing: [3:10:2])
    print("7. Index 3 till 9 (varannat element):", data[3:10:2])

    # 8. Printa ut listan baklänges med slicing
    print("8. Listan baklänges (slicing):", data[::-1])

    # 9. Printa ut listan med en for-loop
    print("9. Utskrift med for-loop:")
    for item in data:
        print(item, end=" ")
    print()

    # 10. Printa ut listan med en while-loop
    print("10. Utskrift med while-loop:")
    i = 0
    while i < len(data):
        print(data[i], end=" ")
        i += 1
    print()

    # 11. Hur lång är listan?
    print("11. Längd på listan:", len(data))

    # 12. Finns 7 i listan?
    print("12. Finns 7 i listan?:", 7 in data)

    # 13. Vilket index har 'katt'?
    print("13. Index för 'katt':", data.index("katt"))

    # 14. Hur många 'katt' finns i listan?
    print("14. Antal 'katt' i listan:", data.count("katt"))

    # 15. Printa listan baklänges med .reverse()
    data.reverse()
    print("15. Listan efter .reverse():", data)

    # 16. Lägg till en 'hund' på slutet med insert
    data.insert(len(data), "hund")
    print("16. Efter .insert() på slutet med 'hund':", data)
    \end{lstlisting}
    \item Problem 6
    \begin{lstlisting}
        import matplotlib.pyplot as plt

        # Funktion med standardvärde step=1
        def generera_x(start, stop, step=1):
            return list(range(start, stop, step))

        # 1. Standardläge (steg 1)
        x_val = generera_x(0, 10)
        y_val = [x**2 for x in x_val]

        # 2. Om användaren anger eget steg (t.ex. steg 2)
        x_anpassad = generera_x(0, 10, step=2)
        y_anpassad = [x**2 for x in x_anpassad]

        # Rita grafen
        plt.plot(x_val, y_val, label="Standard (steg 1)")
        plt.plot(x_anpassad, y_anpassad, label="Anpassad (steg 2)")

        plt.title("Graf med range() och Matplotlib")
        plt.xlabel("x")
        plt.ylabel("y")
        plt.legend()
        plt.grid(True)
        plt.show()
    \end{lstlisting}
    \item Problem 7
    \begin{lstlisting}
        mening = "jag grät när jag inte fick mina björnbär"
        forbjudna_bokstaver = "åäö"

        # 1. Undersök om meningen innehåller någon av bokstäverna
        har_forbjuden = False
        for bokstav in forbjudna_bokstaver:
            if bokstav in mening:
                har_forbjuden = True

        print("Innehåller meningen förbjudna bokstäver?:", har_forbjuden)

        # 2. Hitta och skriv ut orden
        ord_lista = mening.split()
        ord_med_forbjudna = []

        for ordet in ord_lista:
            # Kolla om ordet innehåller å, ä eller ö
            if (
                ("å" in ordet)
                or ("ä" in ordet)
                or ("ö" in ordet)
            ):
                ord_med_forbjudna.append(ordet)

        print("Ord som innehåller förbjudna bokstäver:", ord_med_forbjudna)
    \end{lstlisting}
    \item Problem 8
    \begin{lstlisting}
        # 1. Skapa en dictionary med nycklar och värden
        lager = {"äpple": 15, "banan": 8, "päron": 12, "apelsin": 20}

        print("Ursprunglig dictionary:", lager)

        # 2. Ta fram ett värde med .get()
        antal_applen = lager.get("äpple")
        print("Antal äpplen i lager:", antal_applen)

        # 3. Uppdatera ett värde
        lager["banan"] = 25
        print("Efter uppdatering av banan:", lager)

        # 4. Loopa över alla par med .items()
        print("\n--- LAGERSTATUS ---")
        for frukt, antal in lager.items():
            print(f"Frukt: {frukt:<8} | Antal i lager: {antal}")
    \end{lstlisting}
    \item Problem 9
    \begin{lstlisting}
        # 1. Läs in kattbrevet med readlines() och with open
        with open("karleksbrev_katt.txt", "r") as fil:
            rader = fil.readlines()

        # 2. Ersätt 'katt' med 'hund' på varje rad och spara i en ny fil
        with open("karleksbrev_hund.txt", "w") as ny_fil:
            for rad in rader:
                ny_rad = rad.replace("katt", "hund").replace("Katt", "Hund")
                ny_fil.write(ny_rad)

        # 3. Läs in den nya filen och printa ut i konsolen
        with open("karleksbrev_hund.txt", "r") as ny_fil:
            innehall = ny_fil.read()
            print(innehall)
    \end{lstlisting}
    \item Problem 10
    \begin{lstlisting}
        # 1. Bygg 3x3-matrisen
        matris = []
        num = 1
        for r in range(3):
            rad = []
            for k in range(3):
                rad.append(num)
                num += 1
            matris.append(rad)

        # 2. Skriv ut matrisen och beräkna summan
        total = 0
        for rad in matris:
            for tal in rad:
                print(tal, end=" ")
                total += tal
            print()

        print("Summa:", total)
    \end{lstlisting}
    \item Problem 11
    \begin{lstlisting}
        # 1. While loop som räknar ner från 5 till 0
        nummer = 5
        while nummer >= 0:
            print(nummer)
            nummer -= 1
    \end{lstlisting}
    \item Problem 12
    \begin{lstlisting}
        def odd_or_even(num):
            if num % 2 == 0:
                print(f"{num} är even")
            else:
                print(f"{num} är odd")

        # Exempelanrop
        odd_or_even(4)  # Output: 4 är even
        odd_or_even(7)  # Output: 7 är odd
    \end{lstlisting}
    \item Problem 13
    \begin{lstlisting}
        def is_longer_than_8(word):
            return len(word) > 8

        # Exempelanrop
        print(is_longer_than_8("programmering"))  # Output: True
        print(is_longer_than_8("kort"))           # Output: False
    \end{lstlisting}
\end{enumerate}

\end{document}